% =====================================================================
% Documentation of the AI-based tools used.
% One row per tool and purpose:
%   \kihilfsmittel{tool}{form of use}{affected parts}{prompts used}
% The table itself is typeset in diplomarbeit.tex via \listekihilfsmittel.
% If no AI tool was used, delete every row -- the table then disappears
% automatically.
% =====================================================================

\kihilfsmittel{Claude Opus (Anthropic)}
  {Ausschließlich zur Prüfung von Rechtschreibung und Grammatik. Inhalte selbst verfasst.}
  {Gesamtes Dokument}
  {\enquote{Überprüfe folgenden Text auf Rechtschreib-, Grammatik- und Beistrichfehler.}}

\kihilfsmittel{Claude Opus (Anthropic)}
  {Ausschließlich zur Übersetzung ins Englische. Inhalt selbst verfasst.}
  {Kapitel \enquote{Abstract}}
  {\enquote{Übersetze folgenden Text ins Englische, ohne den Inhalt zu verändern.}}

\kihilfsmittel{Kimi}
  {Codevervollständigung während der Implementierung. Ergebnisse geprüft und angepasst.}
  {Praktische Realisierung}
  {Inline-Vorschläge, keine expliziten Prompts}

\kihilfsmittel {Gemini / Claude}
  {Stilistische Überarbeitung (Lektorat) und Anpassung vollständig selbst verfasster Textentwürfe an den akademischen Schreibstil.}
  {Gesamtes Dokument}
  {Überarbeite meinen folgenden Textentwurf hinsichtlich Satzbau und akademischem Ausdruck (Passivform). Verändere keine inhaltlichen Fakten und füge keine neuen Informationen hinzu.}

\kihilfsmittel{Codex (OpenAI), 29.10.2026}
  {Erstellung und Überarbeitung von Textentwürfen;
   Abgleich mit Programmcode und Unterstützung bei der Kalibrierungsprüfung.}
  {Teilbereich E. Akpinar: technische Grundlagen und praktische Realisierung.}
  {\cite{akpinar-codex-20261029a,akpinar-codex-20261029b,akpinar-codex-20261029c}}
